\documentclass[12pt]{article}
\usepackage[a4paper,margin=2.5cm]{geometry}
\usepackage{amsmath,amssymb}
\usepackage[indonesian]{babel}
\usepackage{enumitem}

\setlength{\parindent}{0pt}
\setlength{\parskip}{6pt}
\setlist{leftmargin=*, itemsep=6pt}

\begin{document}

\section*{1. Perbedaan Deret Taylor dan Deret Maclaurin}

Deret Taylor fungsi $f(x)$ di sekitar titik $x=a$:
\[
  f(x) = \sum_{n=0}^{\infty} \frac{f^{(n)}(a)}{n!}\,(x-a)^n .
\]
Deret Maclaurin adalah deret Taylor dengan $a=0$:
\[
  f(x) = \sum_{n=0}^{\infty} \frac{f^{(n)}(0)}{n!}\,x^n .
\]
Jadi perbedaannya hanya pada titik pusat ekspansi. Contoh deret Maclaurin:
\[
  e^x = 1 + x + \frac{x^2}{2!} + \frac{x^3}{3!} + \cdots
\]
Contoh deret Taylor di $a=1$ (karena $\ln x$ tidak terdefinisi di $x=0$):
\[
  \ln x = (x-1) - \frac{(x-1)^2}{2} + \frac{(x-1)^3}{3} - \cdots
\]

\section*{2. Syarat Kekonvergenan Deret}

Deret $\sum a_n$ konvergen jika jumlah parsialnya
\[
  S_n = a_1 + a_2 + \cdots + a_n
\]
menuju bilangan berhingga $S$ saat $n \to \infty$.

Deret konvergen hanya jika
\[
  \lim_{n\to\infty} a_n = 0 .
\]
Selain syarat di atas, ada cukup banyak uji yang dapat digunakan untuk
menentukan kekonvergenan suatu deret. Contohnya deret harmonik
$\sum_{n=1}^{\infty} \frac{1}{n}$: sukunya menuju $0$, tetapi deretnya divergen.

\subsection*{Uji yang umum dipakai}

\begin{itemize}
  \item Deret geometri $\sum a r^n$ konvergen jika $|r| < 1$, dengan jumlah
  $\frac{a}{1-r}$.

  Contoh:
  \[
    \sum_{n=0}^{\infty} \left(\frac{1}{2}\right)^{n} = 2 .
  \]

  \item Deret-$p$ $\sum \frac{1}{n^p}$ konvergen jika $p > 1$.

  Contoh:
  \[
    \sum_{n=1}^{\infty} \frac{1}{n^2} = \frac{\pi^2}{6} .
  \]

  \item Uji rasio: jika
  \[
    L = \lim_{n\to\infty} \left| \frac{a_{n+1}}{a_n} \right| < 1,
  \]
  deret konvergen.

  Contoh: untuk $\sum \frac{1}{n!}$,
  \[
    L = \lim_{n\to\infty} \frac{1}{n+1} = 0 < 1,
  \]
  jadi konvergen.

  \item Deret berganti tanda (Leibniz): $\sum (-1)^{n+1} b_n$ konvergen jika
  $b_n$ turun dan $\lim b_n = 0$.

  Contoh:
  \[
    1 - \frac{1}{2} + \frac{1}{3} - \frac{1}{4} + \cdots = \ln 2 .
  \]
\end{itemize}

\end{document}
